%% chapters/00-abstract-en.tex
%% Abstract Bahasa Inggris — content equivalen dengan abstrak ID
%% Aturan: lihat .cursor/rules/12-abstrak-front-matter.mdc

\begin{abstract*}

Remaining Useful Life (RUL) prediction for rolling element bearings
underpins modern predictive maintenance strategies. State-of-the-art
methods based on Long Short-Term Memory (LSTM) networks and
Transformers face a trade-off between long-sequence accuracy,
computational cost, and the interpretability of model decisions with
respect to the underlying physical degradation phenomena of bearings.
This dissertation proposes \emph{Mamba-xLSTM-Net}, a hybrid
architecture that integrates a selective state-space model (Mamba) for
long-range dynamics with an xLSTM module equipped with matrix memory
for precise recurrent state tracking. The architecture is augmented by
a post-hoc interpretability module based on a Top-$k$ Sparse
Autoencoder (SAE) whose latent features are mapped empirically to the
characteristic bearing frequencies (Ball Pass Frequency Outer race,
Ball Pass Frequency Inner race, Ball Spin Frequency, and Fundamental
Train Frequency).

Experiments were conducted on two public datasets, PHM2012 and
XJTU-SY, with baselines including LSTM, Transformer, Informer, vanilla
Mamba, and vanilla xLSTM. Each configuration was trained with five
random seeds and evaluated using Root Mean Square Error (RMSE), Mean
Absolute Error (MAE), and the PHM2012 scoring function. The proposed
method achieved an average RMSE reduction of XX\% and an MAE reduction
of XX\% relative to the strongest baseline, with the largest gains
under high-load operating conditions. Ablation studies confirmed that
the Mamba and xLSTM components contribute additively: removing either
module led to a statistically significant performance drop. The
interpretability analysis revealed that the SAE latent features most
frequently activated during late-life regimes correspond strongly with
spectral components near the theoretical BPFO, BPFI, and BSF,
providing empirical evidence that the model learns representations
consistent with classical bearing vibration diagnostics.

The main contributions of this dissertation are: (i) the
Mamba-xLSTM-Net hybrid architecture that combines the strengths of
selective state-space models and matrix-memory recurrent networks;
(ii) an interpretability methodology that links SAE latent features to
the physical concept of bearing characteristic frequencies; and (iii)
empirical evidence of significant performance improvement on two
widely used benchmarks. The results pave the way for deep
learning-based prognostic systems that are simultaneously more
accurate and more trustworthy for maintenance engineers in industry.

\vspace{0.5cm}
\noindent
\textbf{Keywords:} remaining useful life, rolling bearing, Mamba,
xLSTM, sparse autoencoder, interpretability, predictive maintenance

\end{abstract*}
